% Generated draft: requires mathematical and source review.
\documentclass[11pt]{article}
\usepackage[margin=25mm]{geometry}
\usepackage{fontspec}
\setmainfont{Latin Modern Roman}
\usepackage{amsmath,amssymb,mathtools,unicode-math}
\setmathfont{Latin Modern Math}
\usepackage{enumitem,needspace,xurl,hyperref}
\setlength{\parindent}{0pt}
\setlength{\parskip}{5pt}
\setlength{\emergencystretch}{3em}
\newcommand{\N}{\mathbb{N}}
\newcommand{\Z}{\mathbb{Z}}
\newcommand{\Q}{\mathbb{Q}}
\newcommand{\R}{\mathbb{R}}
\newcommand{\C}{\mathbb{C}}
\newcommand{\F}{\mathbb{F}}
\newcommand{\A}{\mathbb{A}}
\newcommand{\PP}{\mathbb{P}}
\newcommand{\E}{\mathbb{E}}
\newcommand{\eps}{\varepsilon}
\newcommand{\dd}{\,\mathrm{d}}
\begin{document}
\section*{UnsolvedMath ID: 701184}
\textbf{Source problem code:} catalogue entry
\subsection{Definitions and mathematical statement}
Let \(p\) be a prime and let \(d,n\geq 1\) be integers. A group is \(d\)-generated if it has a generating set of size at most \(d\). A \(p\)-group means a group in which every element has finite order a power of \(p\). A group \(G\) is nilpotent if its lower central series, defined by \(\gamma_1(G)=G\) and \(\gamma_{i+1}(G)=[\gamma_i(G),G]\), reaches the trivial subgroup. Here \([A,B]\) denotes the subgroup generated by the commutators \(a^{-1}b^{-1}ab\), with \(a\in A\) and \(b\in B\).

For \(\varphi\in\operatorname{Aut}(G)\), write \(x^{\varphi^i}=\varphi^i(x)\). Following the source, a \emph{splitting automorphism of order \(m\)} is an automorphism satisfying
\[
\varphi^m=\operatorname{id}_G,
\qquad
x\,x^{\varphi}x^{\varphi^2}\cdots x^{\varphi^{m-1}}=1
\quad\text{for every }x\in G.
\]
The product is taken in the displayed order. This terminology does not require that the least positive order of \(\varphi\) equal \(m\).

Define the derived series by \(G^{(0)}=G\) and \(G^{(r+1)}=[G^{(r)},G^{(r)}]\). The derived length \(\operatorname{dl}(G)\) is the least nonnegative integer \(r\) such that \(G^{(r)}=1\).

\textbf{(a)} Does there exist a function
\[
F:\mathbb{Z}_{\geq1}\times\{p:p\text{ is prime}\}\times\mathbb{Z}_{\geq1}
\longrightarrow\mathbb{Z}_{\geq0}
\]
such that, for every \(d\geq1\), every prime \(p\), every \(n\geq1\), and every \(d\)-generated nilpotent \(p\)-group \(G\) admitting a splitting automorphism of order \(p^n\),
\[
\operatorname{dl}(G)\leq F(d,p,n)?
\]
In particular, the bound must be independent of the nilpotency class and of the particular group and automorphism. The source reports that the answer is affirmative for \(n=1\), referring to Archive, Problem 7.53; this reported result is not independently verified here.

\textbf{(b)} The source separately asks the same question for \(p^n=4\). Under the hypotheses of part (a), this means \(p=2\) and \(n=2\): does there exist a function \(B:\mathbb{Z}_{\geq1}\to\mathbb{Z}_{\geq0}\) such that every \(d\)-generated nilpotent \(2\)-group admitting an automorphism \(\varphi\) satisfying
\[
\varphi^4=\operatorname{id}_G,
\qquad
x\,x^{\varphi}x^{\varphi^2}x^{\varphi^3}=1
\quad\text{for every }x\in G
\]
has \(\operatorname{dl}(G)\leq B(d)\)? No different hypotheses for part (b) are specified in the supplied record.
\subsection{Short English statement}
E. I. Khukhro asks whether the derived length of a nilpotent \(p\)-group generated by at most \(d\) elements and admitting a splitting automorphism of order \(p^n\) is bounded solely in terms of \(d,p,n\). The source separately highlights order \(4\), and reports an affirmative answer for \(n=1\).
\subsection{Sources}
\begin{itemize}
\item Supplied local UnsolvedMath registry record: permanent ID \(701184\), exact source problem code \texttt{catalogue entry}, title \emph{Kourovka Notebook Problem 15.96}; upstream revision \texttt{d371ee917ac9647fbb06cd6fc529deec6f3dda82}. The record was created and updated on \texttt{2026-10-16T14:00:00Z}.
\item \emph{The Kourovka Notebook (Unsolved Problems in Group Theory)}, 15th issue (2002), Problem 15.96; proposed by E. I. Khukhro in 2002. The supplied background also identifies the 21st issue edition (2026), arXiv:\texttt{1401.0300v48}. The collection description names E. I. Khukhro and V. D. Mazurov as editors.
\item Archive, Problem 7.53: the supplied statement cites this for its reported affirmative answer when \(n=1\). No fuller citation for that entry is supplied.
\item The registry labels this problem \texttt{open}. This is a supplied source label, not an independently verified assessment of the literature. The cited editions and archive entry have not been independently consulted for this draft.
\end{itemize}
\end{document}
